%============================================================
\chapter{前受型と成果連動型の実証：族2と族5}

%============================================================

\section{本章で検証する命題}

第\ref{part:catalog}部で列挙した類型のうち、族2（継続アクセス）と族5（成果・状態連動）を対象に、
命題\ref{prop:breakage}（乖離の発生条件）と第\ref{sec:info}章の可測性条件を検証する。
第\ref{part:method}部の方針に従い、調査に先立って予測を確定し、
その予測が支持されたか否かを、外れた部分を含めてそのまま記載する。

事前に確定した予測は次の二つである。

\begin{quote}
\textbf{予測A（族2）}：取引頻度 $\nu$ が低い商品ほど、権利と行使の乖離が大きい。
具体的には、日次利用される決済手段より、贈答用の商品券のほうが乖離が大きい。

\textbf{予測B（族5）}：成果の事後検証が難しい職種ほど、返戻金制度が強い
（返戻率が高い、または対象期間が長い）。
\end{quote}

結論を先に述べる。予測Aは水準について支持されたが、
測定できたのは $\phi_{\mathrm{cog}}$ ではなく $\kappa$ であった。
予測Bは検証できなかった。必要な変数が集計統計に存在しない。


\begin{center}
\begin{tabularx}{\textwidth}{@{}llX@{}}
\toprule
\multicolumn{3}{@{}l}{\textbf{本章で扱う理論量}}\\
\midrule
$D,\ P,\ \kappa_C$ & 式\eqref{eq:kappa} & 発行額・回収額・未使用残高として直接観測 \\
$\CCC = W/r$ & 式\eqref{eq:ccc} & 媒体別・業種別に算出。15--1{,}904 日 \\
$\dbar-\delta$ & 第\ref{sec:breakage}節 & 有効期限到来分。上界のみ \\
$\nu$ & 式\eqref{eq:decay} & 購入事由（贈答比率）で代理 \\
$\mathcal{G}$ & 第\ref{sec:info}章 & 掲載項目集合として族5で観測 \\
$y$ & 第\ref{sec:info}章 & 6か月以内離職者数 \\
$b$ & 例\ref{ex:bstar} & 返戻金制度の有無。率と期間は不統一 \\
手数料実績率 & 第\ref{sec:predX-result}節 & 価格側の代理変数。連続量 \\
\bottomrule
\end{tabularx}
\captionof{table}{本章（族2・族5）で扱う理論量と対応する実測値。}
\label{tab:family2-5-theory-quantities}
\end{center}

\section{族2：前払式支払手段}

\subsection{データ源と母集団}

一般社団法人日本資金決済業協会「第27回発行事業実態調査統計」（令和6年度版）を用いた。
第\ref{sec:frame}節の方針どおり、上場・非上場を問わない全数フレームである。

\begin{center}
\begin{tabularx}{\textwidth}{@{}lX@{}}
\toprule
項目 & 内容 \\
\midrule
母集団 & 令和6年度末の登録・届出発行者 2,034者（第三者型829者・自家型1,205者） \\
回答 & 812者、回答率 39.9\%。ただし発行額ベースでは全発行者の約9割をカバー \\
規模 & 発行額 28.2兆円、未使用残高 2.90兆円 \\
制約 & 協会は「調査ごとの回答者が異なるため、計数には連続性はありません」と明記 \\
\bottomrule
\end{tabularx}
\captionof{table}{族2（前払式支払手段）のデータ源と母集団（令和6年度）。}
\label{tab:family2-data-source}
\end{center}

\subsection{滞留期間の測定}

本節の量を理論の記号に対応させる。発行額が累積決済 $P$、回収額が累積履行 $D$、
未使用残高が式\eqref{eq:kappa}の $-\kappa_C$ にあたる。令和6年度の値は
\begin{equation}
P = 28{,}150{,}350,\quad D = 28{,}241{,}983,\quad \kappa_C = -2{,}895{,}327
\qquad\text{（百万円）}
\end{equation}
であり、$\kappa_C<0$ すなわち顧客が発行者に与信している。
運転資本を $W=-\kappa_C$、売上速度を $r=P$ とすれば式\eqref{eq:ccc}より
\[
\CCC = \frac{W}{r} = \frac{2{,}895{,}327}{28{,}150{,}350}\times 365 = 37.5\ \text{日}
\]
を得る。これを日数に換算したものを滞留日数と呼ぶ。業種別の結果は表\ref{tab:family2-industry-durations}のとおりである。

\begin{center}
\begin{tabular}{@{}lrr@{}}
\toprule
業種 & 発行額（百万円） & 滞留日数 \\
\midrule
ホテル・旅館業 & 652 & 2,111 \\
旅行業 & 42,603 & 1,948 \\
百貨店 & 72,890 & 1,223 \\
不動産業 & 1,230 & 650 \\
通信業 & 237,299 & 463 \\
協同組合・商工会議所等 & 187,364 & 140 \\
スポーツ・レジャー & 12,361 & 115 \\
外食業 & 142,327 & 107 \\
クレジット・割賦販売業 & 4,301,983 & 59 \\
運輸業 & 3,123,308 & 26 \\
スーパー & 2,198,404 & 17 \\
\midrule
全体 & 28,150,350 & 38 \\
\bottomrule
\end{tabular}
\captionof{table}{業種別の発行額と滞留日数（族2、令和6年度）。}
\label{tab:family2-industry-durations}
\end{center}

順序は予測Aと整合する。運輸業（交通系IC、日次利用）とスーパー（日常買物）が
20日前後であるのに対し、宿泊券・旅行券・百貨店商品券は1,200日から2,100日である。

補強材料として、式\eqref{eq:decay}の取引頻度 $\nu$ の代理変数となる購入事由の集計がある。
紙型は購入者自身が使うものが41.5\%、贈答用が30.1\%であるのに対し、
IC型は自身が使うものが95.5\%、贈答用が0.0\%である。
贈答では $\bar\delta$ を保有する主体と $\delta$ を実行する主体が分離するため、
$\nu$ が構造的に低くなる。

\subsection{対抗仮説の分離}

有力な対抗仮説がある。紙商品券市場の縮小により発行額（フロー）が減少し、
未使用残高（ストック）が過去の蓄積として残れば、比率は機械的に上昇する。
これは第\ref{sec:apc}節の時点効果にほかならない。

業種別の断面では分離できない。しかし媒体別のパネルには、
\textbf{成長している媒体と縮小している媒体が同時に含まれる}ため、識別が可能になる。

\begin{center}
\begin{tabular}{@{}lrrrrrr@{}}
\toprule
媒体 & R2 & R3 & R4 & R5 & R6 & R2$\to$R6 発行額 \\
\midrule
磁気型 & 1,307 & 1,630 & 1,338 & 1,916 & 1,904 & $-32.2\%$ \\
紙型 & 775 & 758 & 668 & 764 & 845 & $-6.0\%$ \\
サーバ型 & 16 & 15 & 17 & 16 & 17 & $+45.3\%$ \\
IC型 & 15 & 14 & 16 & 15 & 15 & $+23.5\%$ \\
\bottomrule
\end{tabular}
\captionof{table}{媒体別の滞留日数の推移と発行額変化率（R2〜R6年度、族2）。}
\label{tab:family2-media-panel}
\end{center}
\begin{center}
\footnotesize（数値は滞留日数。最右列のみ発行額の変化率）
\end{center}

判定は二つに分かれる。

\begin{enumerate}[label=(\arabic*)]
\item \textbf{水準差は市場縮小では説明できない。}
紙型は5年間で発行額が6.0\%しか減っていないにもかかわらず、
令和2年度の時点ですでに775日の滞留がある。IC型の15日に対して約50倍である。
縮小が進行していない起点で差が存在する以上、水準は構造に由来する。
また成長している二媒体（IC型・サーバ型）の滞留は5年間ほぼ不動であり、
機械的な比率上昇は起きていない。

\item \textbf{変動の一部は市場縮小で説明できる。}
磁気型は発行額が32.2\%減少し、滞留は1,307日から1,904日へ上昇した。
紙型も令和4年度の668日から845日へ戻している。ここは対抗仮説が効いている。
\end{enumerate}

\begin{figure}[htbp]
\centering
\insertfig{media-panel}
\caption{媒体別の滞留日数（対数目盛）。二桁の水準差が起点から存在する。}
\label{fig:media-panel}
\end{figure}

すなわち\textbf{水準は $\nu$、変動の一部は時点効果}という分離になる。
第\ref{sec:apc}節で一般には識別不能と述べた問題に対し、
成長セグメントが同時に存在したことで部分的な識別ができた。
これは一般解ではなく、この題材に固有の幸運である。

補助的な観測として、使用期限を設定していない比率は
紙型71.2\%、磁気型52.0\%、IC型37.9\%、サーバ型（リアル店舗）19.0\%であり、
磁気型を除けば滞留日数と順序が一致する。
ただし因果の向きは特定できない。期限を設定しないから滞留するのか、
贈答用のように滞留する商品だから期限を設定しないのか、
おそらく同じ $\nu$ が双方を決めているが、このデータでは決着しない。

\subsection{検証できなかったこと}

\begin{proposition}[測定対象の限定]
本節で測定できたのは信用ポジション $\kappa$ の大きさであって、
式\eqref{eq:cog}の認知余剰 $\phi_{\mathrm{cog}}$ ではない。
\end{proposition}

滞留が長いことは、権利が最終的に行使されないことを含意しない。
$\phi_{\mathrm{cog}}$ に対応する数値は有効期限到来等による回収額 20,634百万円のみであり、
未使用残高比 0.71\%、発行額比 0.073\% にすぎない。
しかも協会はこの金額に払戻しによる回収額も含まれているものと考えられると注記しており、
純粋な非行使額ではない。

したがって族2 の結論は次に限定される。

\begin{quote}
前受型 $\Phi$ において、取引頻度 $\nu$ が低い商品ほど $|\kappa|$ が大きく、
その差は二桁に及ぶ。ただしこの滞留が認知余剰に転化するかは、公開データでは判定できない。
\end{quote}

第\ref{sec:growth}章（$\kappa$ と滞留）は支持されたが、
第\ref{sec:breakage}節（認知余剰）は未検証のまま残る。前者は後者の必要条件にすぎない。

\subsection{個社データの不在}

協会は、財務局等行政機関職員には法令上守秘義務が課せられており、
発行者から報告された事項を第三者である協会に情報提供することはできないと明記している。
金融庁は登録・届出者の一覧を公表しているが、そこに未使用残高は含まれない。
加えて、発行者自身から未回収率（期限切れ等）の状況も開示してほしいという要望が出ており、
すなわち退蔵率は現在集計されていない。

結果として、業種内・媒体内のばらつきは一切観測できない。見えているのは平均だけである。

\section{族5：有料職業紹介}

\subsection{データ源と母集団}

厚生労働省「令和6年度職業紹介事業報告書の集計結果（速報）」を用いた。
職業安定法に基づく全事業者の悉皆報告であり、族2 と同様に上場バイアスを免れる。

\begin{center}
\begin{tabularx}{\textwidth}{@{}lX@{}}
\toprule
項目 & 内容 \\
\midrule
母集団 & 事業報告を提出した有料職業紹介事業所 30,561。うち就職実績のあった事業所 13,946 \\
規模 & 常用就職 888,993件（有料）、手数料収入 約9,835億円 \\
単価 & 常用就職1件あたり手数料 約103万円 \\
\bottomrule
\end{tabularx}
\captionof{table}{族5（有料職業紹介）のデータ源と母集団（令和6年度）。}
\label{tab:family5-data-source}
\end{center}

\subsection{契約レジームの選択}

職業安定法は手数料徴収方法として上限制と届出制を用意している。
上限制は賃金額の一定率に固定された公式ベースの契約、
届出制は事業者が自由に設定して届け出る契約である。

\begin{center}
\begin{tabular}{@{}lrr@{}}
\toprule
手数料の種類 & 金額（千円） & 構成比 \\
\midrule
届出制手数料 & 980,821,781 & 99.73\% \\
上限制手数料 & 1,804,568 & 0.18\% \\
その他（受付等） & 827,808 & 0.08\% \\
\midrule
合計 & 983,454,157 & 100\% \\
\bottomrule
\end{tabular}
\captionof{table}{手数料徴収方法別の内訳（族5、令和6年度）。}
\label{tab:family5-fee-type-breakdown}
\end{center}

上限制はほぼ消滅している。

\begin{remark}[この事実は可測性の証拠ではない]
上限制は賃金額の約11\%（6か月分）に固定されており、年収換算では5\%程度にしかならない。
届出制の実勢である年収の30\%前後とは水準が一桁異なる。
すなわち市場が公式ベースの契約を拒否したのではなく、公式が定める価格が低すぎるだけである。
これは制度層が価格を規定した結果として選択肢が使われなくなった現象であり、
第\ref{sec:info}章の可測性とは無関係である。
上限制の消滅は価格水準の帰結であり、可測性仮説の材料にはならない。
\end{remark}

\subsection{職種別の1件あたり手数料}

手数料には臨時日雇分が混在している。全国計では常用就職に係る手数料が
9,136億円と全体の約93\%を占めるが、職種別の内訳では両者が分離されていない。
そこで、常用就職1件あたりの臨時日雇就職延数（人日）を混入度と定義し、
これが1未満の職種に限って算出した。

\begin{center}
\begin{tabular}{@{}lrrr@{}}
\toprule
職種 & 就職件数 & 混入度 & 万円/件 \\
\midrule
法人・団体役員 & 1,392 & 0.42 & 437 \\
法人・団体管理職員 & 16,007 & 0.16 & 288 \\
経営・金融・保険の専門的職業 & 11,346 & 0.09 & 268 \\
法務の職業 & 2,217 & 0.05 & 225 \\
研究者 & 3,167 & 0.27 & 188 \\
総務・人事・企画事務 & 25,716 & 0.38 & 183 \\
開発技術者 & 19,339 & 0.40 & 165 \\
情報処理・通信技術者（ソフト開発） & 39,780 & 0.20 & 150 \\
営業の職業 & 89,563 & 0.28 & 142 \\
建築・土木・測量技術者 & 26,230 & 0.12 & 139 \\
デザイナー & 3,969 & 0.38 & 116 \\
\bottomrule
\end{tabular}
\captionof{table}{職種別の1件あたり手数料（混入度1未満の職種、族5）。}
\label{tab:family5-fee-per-placement}
\end{center}

序列は職位の高さに対応しており、手数料率がほぼ一定であれば年収の反映にすぎない。
\textbf{年収の分母がないため、可測性仮説に対する識別力がない。}
役員の手数料が高いのが、成果の観測しにくさによるのか単に年収が高いからなのか、
このデータでは区別できない。

\subsection{予測Bが検証できない理由}

返戻金制度のデータは本集計に存在しない。個社ごとに人材サービス総合サイトに
掲載されているのみである。したがって予測Bは未検証である。

さらに深刻な問題がある。

\begin{remark}[データの穴が関心と一致する]
医療・介護・保育は、返戻金制度を含む規制が集中的に投入された分野である。
制度層が介入したという事実自体が、この分野で可測性の問題が深刻だと
当局が判断した証拠になる。ところがこれらの職種は臨時日雇（スポット勤務）の
混入が著しく、集計値から1件あたり手数料を算出できない。
混入度は医師35、看護師・准看護師12、施設介護5、保育士4である。

すなわち\textbf{理論的関心が最も高い領域が、そのままデータの穴になっている}。
これは偶然ではない可能性がある。スポット勤務が発達する分野は、
雇用関係が短期化しやすく、まさに成果の事後検証が難しい分野だからである。
\end{remark}

\section{二族の比較と方法論的含意}

\begin{center}
\begin{tabularx}{\textwidth}{@{}lXX@{}}
\toprule
 & 族2 & 族5 \\
\midrule
予測 & 水準について支持 & 未検証 \\
集計層で分かること & $\kappa$ の水準差が二桁 & ほぼ無し \\
対抗仮説 & 部分的に分離できた & 未着手 \\
個社データ & 不可（守秘義務） & 可（人材サービス総合サイト） \\
残る課題 & $\phi_{\mathrm{cog}}$ の測定 & 抽出設計と個社データ取得 \\
\bottomrule
\end{tabularx}
\captionof{table}{族2と族5の比較：検証状況・個社データ可否・残る課題。}
\label{tab:family2-5-comparison}
\end{center}

第\ref{part:catalog}部末尾で、族5 のほうが条件が良いと見込んだ。
これは個社データに到達できればという条件つきの見込みであり、
集計層のみを見た現時点では族2 より弱い結論しか得られていない。
\textbf{フレームの網羅性と変数の充実は独立であり、両方を確認しないと調査の見込みは立たない。}

もう一つ、両族に共通する構造が観察された。

\begin{proposition}[集計層と個社層のトレードオフ]
本稿が扱う二つの領域では、規制に基づく全数フレームは構造変数（$\kappa$、契約形態）を
集計値としてのみ提供し、余剰変数（$\phi_{\mathrm{cog}}$、$\phi_{\mathrm{barg}}$）は
選択された部分標本でしか観測できない。
\end{proposition}

族2 で $\phi_{\mathrm{cog}}$ に踏み込むには、退蔵益を個別開示している上場企業の
有価証券報告書に降りる必要があるが、それは第\ref{sec:frame}節で述べた
0.1\%の標本に戻ることを意味し、第\ref{sec:disclosure}節の開示選択バイアスが復活する。
\textbf{構造は全数で見えるが、余剰は選択された標本でしか見えない}という制約が、
少なくともこの二領域では成立している。

\section{事前登録の効果}

第\ref{part:method}部で事前登録を要求したことの効果を記録しておく。

族5 において、上限制の消滅という事実は可測性仮説を支持する材料として
解釈することが可能であった。実際には価格水準の問題であり、無関係である。
予測を事前に確定していなければ、この事実を後付けで理論の傍証として
採用していた可能性が高い。第\ref{sec:narrative}節で述べた事後的物語化は、
他者の記述だけでなく自らの解釈にも生じる。

同様に族2 では、予測が「乖離が大きい」であったのに対し測定できたのは「滞留が長い」であり、
両者の差を明示できたのは予測を先に文章として固定していたからである。

\section{枠組みの改訂点}

以上を踏まえ、第\ref{part:theory}部に対する改訂を三点記録する。

\begin{enumerate}[label=(\arabic*)]
\item \textbf{$\kappa$ と $\phi_{\mathrm{cog}}$ の測定可能性は大きく異なる。}
前者は規制開示から全数で観測できるが、後者は原則として観測できない。
式\eqref{eq:cog}は定義としては明確だが、実証的には空である可能性を認めるべきである。

\item \textbf{制度層による価格規定は、$\Phi$ の選択肢を実質的に消滅させうる。}
族7-3（公的支払）で述べた性質は、公定価格そのものだけでなく、
選択可能な契約形態の上限を通じても作用する。
上限制手数料はこの例であり、注\ref{rem:institution-modes}の第二の形式にあたる。

\item \textbf{データの欠落は無作為ではない。}
族5 で観察されたとおり、理論的関心の高い領域ほど測定が難しいという相関がありうる。
第\ref{part:method}部で扱った選択バイアスは標本の選択に関するものであったが、
\emph{変数}の選択にも同型の問題が存在する。これは新たに記録すべき第七のバイアスである。
\end{enumerate}

\section{次段階の課題}

族5 の個社データ取得に進む。全許可事業者が3万を超えるため全数は現実的でなく、
職種による層化抽出の設計が必要である。設計にあたり、以下を事前に確定する。

\begin{enumerate}[label=(\arabic*)]
\item 層化の基準となる職種区分と、各層の抽出数
\item 混入度が高い職種を含めるか否か、含める場合の扱い
\item 返戻金制度の「強さ」の操作的定義（返戻率、対象期間、またはその組合せ）
\item 年収データの入手経路（可測性仮説の識別に不可欠）
\end{enumerate}

また、選択バイアスの経路として次を記録しておく。
返戻金制度を設けていない事業者は、離職者数を返戻実績で代替できないため、
離職率の開示コストが相対的に高くなる。
すなわち\textbf{開示の有無が説明変数と相関する可能性}があり、
これは第\ref{sec:disclosure}節と同型の問題である。
